
\documentclass{article}
\usepackage{colortbl}
\usepackage{makecell}
\usepackage{multirow}
\usepackage{supertabular}

\begin{document}

\newcounter{utterance}

\centering \large Interaction Transcript for game `cladder', experiment `full\_v1.5\_default', episode 188 with selfrefine{-}llama{-}3.1{-}8b{-}instruct.
\vspace{24pt}

{ \footnotesize  \setcounter{utterance}{1}
\setlength{\tabcolsep}{0pt}
\begin{supertabular}{c@{$\;$}|p{.15\linewidth}@{}p{.15\linewidth}p{.15\linewidth}p{.15\linewidth}p{.15\linewidth}p{.15\linewidth}}
    \# & \multicolumn{2}{c}{Player} && \multicolumn{2}{c}{Game Master} \\
    \hline

    \theutterance \stepcounter{utterance}  
    & & & \multicolumn{4}{p{0.6\linewidth}}{
        \cellcolor[rgb]{0.9,0.9,0.9}{
            \makecell[{{p{\linewidth}}}]{
                \texttt{\tiny{[P1$\langle$GM]}}
                \texttt{You are an expert in causal inference. The following question is not a typical commonsense query, but rather a meticulously designed question created by a professor specializing in causal inference, intended to assess the students' mastery of the course content.} \\
\texttt{Imagine a self{-}contained, hypothetical world with only the following conditions, and without any unmentioned factors or causal relationships: The captain has a direct effect on the private and the corporal. The corporal has a direct effect on prisoner. The private has a direct effect on prisoner. Method 1: We look directly at how the captain correlates with prisoner in general. Method 2: We look at this correlation case by case according to the private. To understand how the captain affects prisoner, is it more correct to use the Method 1 than Method 2?} \\
\texttt{Start your answer with 'Yes' or 'No', followed by additional reasoning or evidence to support your explanation.} \\
            }
        }
    }
     \\ \\

    \theutterance \stepcounter{utterance}  
    & & \multicolumn{4}{p{0.6\linewidth}}{
        \cellcolor[rgb]{0.95,0.95,0.95}{
            \makecell[{{p{\linewidth}}}]{
                \texttt{\tiny{[P1$\leftrightarrow$P]}}
                \texttt{{[}DRAFT{]}} \\
\texttt{No.} \\
\\ 
\texttt{In this hypothetical world, the captain has a direct effect on both the private and the corporal, who in turn have direct effects on the prisoner. This creates a situation where the captain's effect on the prisoner is mediated by both the private and the corporal.} \\
\\ 
\texttt{Method 1, which looks directly at the correlation between the captain and the prisoner in general, is likely to confound the direct effect of the captain on the prisoner with the indirect effects mediated by the private and the corporal. This is because the correlation between the captain and the prisoner will capture both the direct effect of the captain on the prisoner and the effects of the captain on the private and corporal, which in turn affect the prisoner.} \\
\\ 
\texttt{On the other hand, Method 2, which looks at the correlation between the captain and the prisoner case by case according to the private, can help to control for the indirect effects mediated by the private. By conditioning on the private, we can isolate the effect of the captain on the prisoner that is not mediated by the private, which is the direct effect of the captain on the prisoner.} \\
\\ 
\texttt{Therefore, Method 2 is more likely to provide a more accurate estimate of the direct effect of the captain on the prisoner, and is thus the more correct approach.} \\
            }
        }
    }
    & \\ \\

    \theutterance \stepcounter{utterance}  
    & & \multicolumn{4}{p{0.6\linewidth}}{
        \cellcolor[rgb]{0.95,0.95,0.95}{
            \makecell[{{p{\linewidth}}}]{
                \texttt{\tiny{[P1$\leftrightarrow$P]}}
                \texttt{{[}SELF{-}FEEDBACK{]}} \\
\texttt{FORMAT: ok} \\
\texttt{RULES: ok} \\
\texttt{STRATEGY: ok} \\
\texttt{REASON: The answer provides a clear and logical explanation of why Method 2 is more correct than Method 1, effectively addressing the question posed by the game master.} \\
            }
        }
    }
    & \\ \\

    \theutterance \stepcounter{utterance}  
    & \multicolumn{4}{p{0.6\linewidth}}{
        \cellcolor[rgb]{0.561,0.737,0.561}{
            \makecell[{{p{\linewidth}}}]{
                \texttt{\tiny{[P1$\rangle$GM]}}
                \texttt{No.} \\
\\ 
\texttt{In this hypothetical world, the captain has a direct effect on both the private and the corporal, who in turn have direct effects on the prisoner. This creates a situation where the captain's effect on the prisoner is mediated by both the private and the corporal.} \\
\\ 
\texttt{Method 1, which looks directly at the correlation between the captain and the prisoner in general, is likely to confound the direct effect of the captain on the prisoner with the indirect effects mediated by the private and the corporal. This is because the correlation between the captain and the prisoner will capture both the direct effect of the captain on the prisoner and the effects of the captain on the private and corporal, which in turn affect the prisoner.} \\
\\ 
\texttt{On the other hand, Method 2, which looks at the correlation between the captain and the prisoner case by case according to the private, can help to control for the indirect effects mediated by the private. By conditioning on the private, we can isolate the effect of the captain on the prisoner that is not mediated by the private, which is the direct effect of the captain on the prisoner.} \\
\\ 
\texttt{Therefore, Method 2 is more likely to provide a more accurate estimate of the direct effect of the captain on the prisoner, and is thus the more correct approach.} \\
            }
        }
    }
    & & \\ \\

    \theutterance \stepcounter{utterance}  
    & & & \multicolumn{2}{p{0.3\linewidth}}{
        \cellcolor[rgb]{0.9,0.9,0.9}{
            \makecell[{{p{\linewidth}}}]{
                \texttt{\tiny{[GM$|$GM]}}
                \texttt{No} \\
            }
        }
    }
    & & \\ \\

    \theutterance \stepcounter{utterance}  
    & & & \multicolumn{2}{p{0.3\linewidth}}{
        \cellcolor[rgb]{0.9,0.9,0.9}{
            \makecell[{{p{\linewidth}}}]{
                \texttt{\tiny{[GM$|$GM]}}
                \texttt{yes} \\
            }
        }
    }
    & & \\ \\

    \theutterance \stepcounter{utterance}  
    & & & \multicolumn{2}{p{0.3\linewidth}}{
        \cellcolor[rgb]{0.9,0.9,0.9}{
            \makecell[{{p{\linewidth}}}]{
                \texttt{\tiny{[GM$|$GM]}}
                \texttt{game\_result = LOSE} \\
            }
        }
    }
    & & \\ \\

\end{supertabular}
}

\end{document}
